% HD-F1-T26. GPT-6 Astra Pro; /root/transport; 22 November 2026 UTC.
\section[HD-F1-T26: the two-tail Gaussian Specht sources are excluded]{HD-F1-T26: the actual permutation quotient excludes both two-tail Gaussian Specht families in every degree}
\label{sec:f1-gaussian-two-tail-exclusion}

In the residual Gaussian alternating source of T24--T25,
none of
\[
 (n-2,2),\quad (n-2,1,1),\quad
 (n-2,2)',\quad (n-2,1,1)'
 \tag{F1T26.1}
\]
is the coefficient partition. This is a uniform exclusion,
proved using the actual degree-$n$ quotient of the original
cover. It does not assume independent scalar and permutation
projections, and it does not infer a determinant from an
abstract representation replacement.

\paragraph{1. Actual incidence and exterior maps.}
Let $M_1=\mathbb Q^n$ with basis $e_a$, and let
$M_2=\mathbb Q[\binom{[n]}2]$ with orthonormal basis
$e_{\{a,b\}}$. The equivariant incidence maps are
\[
 J(e_a)=\sum_{b\ne a}e_{\{a,b\}},\qquad
 J^*(e_{\{a,b\}})=e_a+e_b,
 \qquad J^*J=(n-2)1+\boldsymbol1\boldsymbol1^t.
 \tag{F1T26.2}
\]
For $n\ge3$ the last operator is invertible. Thus
\[
 M_2=J(M_1)\mathbin{\perp}W_2,\qquad W_2=\ker J^*,
 \qquad\dim W_2=n(n-3)/2.
 \tag{F1T26.3}
\]
Here $W_2$ is the actual rational Specht module of shape
$(n-2,2)$. To verify the identification, identify its row
tabloids with the unordered pair in the second row. The
column-antisymmetrizer on two columns with four different
labels produces
\[
           e_{\{c,d\}}-e_{\{a,d\}}
                    -e_{\{b,c\}}+e_{\{a,b\}},
\]
a nonzero vector in $\ker J^*$. Its translates are the
image of the actual Young-symmetrizer module. The usual
left-ideal convention $\mathbb Q[S_n]ab$ maps to the
tabloid convention $\mathbb Q[S_n]ba$ by right
multiplication by $a$. More explicitly, if
$(ab)^2=\kappa ab$, then the group-ring adjoint gives
$(ba)^2=\kappa ba$, and right multiplication by
$\kappa^{-1}b$ is its inverse. The nonzero tabloid
map from this irreducible module and the Frobenius
dimension $n(n-3)/2$ therefore identify its image with
all of $\ker J^*$. This provides the equivariant map,
not merely a character equality.

For the other shape let $U=\ker(\sum:M_1\to\mathbb Q)$.
The actual orthogonal splitting
$M_1=\mathbb Q\boldsymbol1\perp U$ induces
\[
 \Lambda^2M_1=(\boldsymbol1\wedge U)
                  \mathbin{\perp}\Lambda^2U,
 \qquad u\longmapsto\boldsymbol1\wedge u.
 \tag{F1T26.4}
\]
The displayed map multiplies the positive coefficient
form by exactly $n$. The rational module
$W_{11}=\Lambda^2U$ is the shape $(n-2,1,1)$.
For an explicit Young-symmetrizer map, send a tabloid
whose two singleton rows are $b,c$ to $e_b\wedge e_c$.
The column antisymmetrizer on the labels $a,b,c$ maps
to
\[
                 2(e_b-e_a)\wedge(e_c-e_a),
\]
which is nonzero and lies in $\Lambda^2U$. These
vectors span $\Lambda^2U$. Its dimension
$(n-1)(n-2)/2$ equals the Frobenius dimension of the
Young-symmetrizer module, proving the identification.

All maps (F1T26.2)--(F1T26.4) remain equivariant after
twisting every term by the same original scalar
character $c$. The exterior operation is performed
on the rational permutation module before that twist;
no replacement of $c$ by $c^2$ occurs. The sign-conjugate
partitions are identified by the explicit sign-twist
map of T22, replacing $c$ by $c\,\operatorname{sgn}\sigma$.
This is still an actual $\mu_4$ character, with its
original correlation to $\sigma$ retained.

\paragraph{2. The exact ordinary two-subset rank.}
Let a permutation $s$ have cycle lengths
$\ell_1,\ldots,\ell_k$, let $d=n-k$ be its permutation
index, and let $e$ and $f$ be its numbers of even cycles
and fixed points. On unordered pairs, two distinct point
cycles give $\gcd(\ell_i,\ell_j)$ orbits, each of length
$\operatorname{lcm}(\ell_i,\ell_j)$. Within a point cycle
of length $\ell$, there are $\lfloor\ell/2\rfloor$
pair orbits. Therefore the number of pair orbits is
\[
 K=\sum_i\lfloor\ell_i/2\rfloor
                    +\sum_{i<j}\gcd(\ell_i,\ell_j),
 \qquad
 a_2(s)=n(n-3)/2-K+k.
 \tag{F1T26.5}
\]
The latter equality follows from the actual splitting
(F1T26.3): fixed spaces are exact in characteristic zero.
Put
\[
 \Delta=\sum_{i<j}
          (\ell_i+\ell_j-2\gcd(\ell_i,\ell_j))\ge0.
\]
Direct substitution into (F1T26.5) gives
\[
                    2a_2(s)-(n-6)d=3d-e+\Delta.
 \tag{F1T26.6}
\]
For $s\ne1$, write $r=k-f$ for its number of nontrivial
point cycles. Then $d\ge r$, $e\le r$, and the pairs
consisting of a fixed point and a nontrivial cycle alone
give $\Delta\ge fd$. Since $n=f+r+d$, it follows that
\[
 3d-e+\Delta-n
       \ge 3d-r+fd-(f+r+d)
       =2(d-r)+f(d-1)\ge0.
 \tag{F1T26.7}
\]
Thus the exact ordinary two-subset constituent satisfies
\[
                   a_2(s)\ge\frac{(n-6)d+n}{2}
                         \qquad(s\ne1).
 \tag{F1T26.8}
\]

\paragraph{3. Every original Gaussian phase.}
Let $z\in\mu_4$ be nontrivial, of order $q=2$ or $4$.
For a permutation orbit of length $L$, the twisted
permutation matrix $zP$ has a one-dimensional fixed
space precisely when $q\mid L$. Let $t$ be the number
of point cycles whose lengths are not divisible by $q$,
and let $B$ be the number of unordered-pair orbits whose
lengths are not divisible by $q$. The splitting
(F1T26.3) gives the exact comparison
\[
          \operatorname{rank}(zR_{W_2}(s)-1)
                         =a_2(s)+B-t.
 \tag{F1T26.9}
\]
Because $q$ is a power of two, the least common multiple
of two lengths not divisible by $q$ is again not divisible
by $q$. Each pair of the $t$ exceptional point cycles
therefore supplies at least one of these $B$ orbits.
Hence $B\ge\binom t2$ and
$B-t\ge\binom t2-t\ge-1$ for every integer $t\ge0$.
Together with (F1T26.8) this proves, for $s\ne1$ and
all $z\in\mu_4$,
\[
 \operatorname{rank}(zR_{W_2}(s)-1)
                    \ge\frac{(n-6)d+n-2}{2}.
 \tag{F1T26.10}
\]

To treat $W_{11}$, realize $\Lambda^2M_1$ on the
oriented pair vectors $e_a\wedge e_b$. An orbit of
unordered pairs closes with sign $+1$, except for the
diameter orbit inside an even point cycle of length
$\ell$: that orbit has length $\ell/2$ and closes
by interchanging the two points, with sign $-1$.
For $z=1$ these diameter orbits lose one fixed vector
each. From (F1T26.4) one obtains exactly
\[
             \operatorname{rank}(R_{W_{11}}(s)-1)
                                      =a_2(s)+e.
 \tag{F1T26.11}
\]

For nontrivial $z$, let $P$ count even point cycles for
which $z^{\ell/2}=-1$, and let $N$ count those for
which $z^{\ell/2}=1$. The fixed-space difference between
$z\Lambda^2M_1$ and $zM_2$ is precisely $P-N$.
The constant line has no fixed vector for this $z$.
Using (F1T26.3)--(F1T26.4) and
$\dim W_{11}=\dim W_2+1$ gives
\[
 \operatorname{rank}(zR_{W_{11}}(s)-1)
                     =a_2(s)+B-t+1-P+N.
 \tag{F1T26.12}
\]
Each of the $P$ diameter orbits is included in $B$,
and it is different from every orbit between two
distinct point cycles. Thus
\[
 B-P+N\ge\binom t2,\qquad
 B-t+1-P+N\ge\frac{(t-1)(t-2)}2\ge0.
\]
Equations (F1T26.8), (F1T26.11) and (F1T26.12) prove
\[
 \operatorname{rank}(zR_{W_{11}}(s)-1)
                     \ge\frac{(n-6)d+n}{2}
                 \qquad(s\ne1,\ z\in\mu_4).
 \tag{F1T26.13}
\]
If $s=1$ and the original coefficient element is
nonidentity, then $z\ne1$ and its rank is the entire
coefficient dimension. For $n\ge7$ that dimension
also satisfies (F1T26.10) or (F1T26.13), with $d=0$.
Thus these bounds cover every actual inertia, including
pure scalar inertias.

\paragraph{4. The actual quotient and the contradiction.}
Let $C\to\mathbb P^1$ be the original full cover and
let $H_1$ be the inverse image in $G$ of the stabilizer
of a letter under the actual projection $\sigma$.
The quotient
\[
                         Y=C/H_1\longrightarrow C/G
 \tag{F1T26.14}
\]
is a connected degree-$n$ map, since $\sigma(G)$ is
$A_n$ or $S_n$. Its local fibre over an original
branch circle is the actual $n$-letter permutation
set, so its ramification index contribution is
$d_j=n-\#\operatorname{cycles}(\sigma(g_j))$.
An original branch with trivial permutation projection
has contribution zero. Riemann--Hurwitz for this same
quotient proves
\[
                         \sum_j d_j=2n-2+2g(Y)\ge2n-2.
 \tag{F1T26.15}
\]
There is no change in the original inertia tuple.

For shape $(n-2,2)$, summing (F1T26.10) and using
$b\ge12$, $n\ge7$, gives
\[
 \begin{split}
 \sum_j a_j
   &\ge(n-6)(n-1)+6(n-2)=n^2-n-6,\\
 2m+6&=n^2-3n+6.
 \end{split}
\]
The first lower bound exceeds the required second
quantity by $2n-12>0$. For shape $(n-2,1,1)$,
(F1T26.13) similarly gives
\[
 \begin{split}
 \sum_j a_j
   &\ge(n-6)(n-1)+6n=n^2-n+6,\\
 2m+6&=n^2-3n+8,
 \end{split}
\]
with positive difference $2n-2$. The same inequalities
apply to the sign-conjugate shapes through the exact
scalar-character change in paragraph 1. This proves
the uniform exclusion (F1T26.1).

\paragraph{5. Original target and framework maps.}
The incidence maps and Young-symmetrizer maps in paragraph
1 identify the original coefficient inside its actual
permutation or oriented-pair module. They induce maps
of the original finite local systems, their cellular
complexes and their parabolic images. The unordered
incidence map also comes from the actual incidence
cover with fibre $\{(a,\{a,b\})\}$, with its two
projections to the one-letter and two-letter quotient
curves; pullback and pushforward give exactly the
displayed incidence matrices. After a scalar twist,
the same map is taken with the original character
coefficient system. The exterior map is the displayed
coefficient map; no exterior power of a Jacobian is
substituted for the original target.

The prescribed Gaussian primitive corner stays in
$\mathbb QG$ throughout. If a different invariant
positive coefficient normalization is used for the
explicit models, the original-corner rank-one operator
and its positive polarization scalar are precisely
(F1T25.14)--(F1T25.15). The contradiction uses invariant
ranks and the actual quotient (F1T26.14), so neither
this scalar nor any corner is discarded. In the
condition-A and Cayley--Hamilton correspondence
programme, these actual maps remove these two
coefficient families from the missing supplier gate.
They add no detecting cycle and make no assertion
about the unresolved determinant classes.
